\documentclass{article}
\usepackage[T1]{fontenc}
\usepackage{lmodern}
\usepackage{graphicx}
\usepackage{amsmath,amssymb}
\usepackage[a4paper,margin=2cm]{geometry}
\usepackage[absolute,overlay]{textpos}
\setlength{\TPHorizModule}{1mm}
\setlength{\TPVertModule}{1mm}
\usepackage[indonesian]{babel}
\usepackage{hyperref}
\setlength{\emergencystretch}{2em}
% unit: Halaman 11

\begin{document}

% === Halaman 11 (python/native) ===

• Contoh di dalam Caesar Cipher, misalkan penyerang tahu cipherteks: 

                        KHOOR ZRUOG

    dan juga tahu plainteks-nya seharusnya mengandung kata:

                       HELLO

    Dari pasangan H→K, penyerang tahu pergeseran (k) = 3.

    Dengan kunci k = 3, semua cipherteks lain bisa didekripsi.

11

\end{document}
